\documentclass[12pt,spanish]{article}
\usepackage[utf8]{inputenc}
\usepackage[T1]{fontenc}
\usepackage{graphicx}       
\usepackage{geometry}       
\geometry{margin=1in}
\usepackage{babel}
\usepackage{enumitem}

\title{\textbf{Análisis Estadístico Integral de Subsidios a Combustibles Fósiles}\\\large Estonia (2010-2024)}
\author{Nestor Uriel Sánchez Cuevas}
\date{\today}

\begin{document}
\maketitle

\section{Introducción}
Este informe técnico automatizado recopila las métricas clave obtenidas a partir del procesamiento de datos de subsidios a combustibles fósiles en Estonia. En total, se analizaron un histórico de 70 registros instrumentales del periodo establecido.

\section{Métricas de Tendencia Central y Dispersión}
A partir de las funciones analíticas programadas en Python, se han obtenido los siguientes estadísticos descriptivos fundamentales para la muestra de subsidios anuales:
\begin{itemize}
    \item \textbf{Media aritmética:} 64.51 M EUR
    \item \textbf{Mediana:} 55.30 M EUR
    \item \textbf{Desviación Estándar:} 31.22 M EUR
\end{itemize}

\section{Visualización de Resultados Integrales}
El procesamiento gráfico consolida en una única visualización el comportamiento de las frecuencias de los subsidios, las relaciones de eventos muestrales mediante diagramas de Venn (Gas Natural y Exención Fiscal) y los modelos de distribución probabilística teórica (Binomial y Poisson):

\begin{figure}[htbp]
    \centering
    \includegraphics[width=0.9\textwidth]{mosaico_subsidios.png}
    \caption{Mosaico de análisis estadístico integral (Proyecto NUSC).}
\end{figure}

\section{Conclusiones y Modelado}
La distribución empírica de los datos permite contrastar el comportamiento histórico real contra las aproximaciones de probabilidad de eventos discretos. Las herramientas de generación dinámica demuestran la viabilidad de reportes automatizados enlazando entornos de ejecución científica con sistemas de composición tipográfica.\\

Fin del reporte.
\end{document}
